%=====================================================
\section{Tier 2.1 revisited: DCB with Box 4.4 arc-length}
\label{sec:dcb_arclength}

The Tier 2.1 DCB benchmark (Turon et al.\ 2006, AS4/PEEK) was previously
reported as \emph{partial-pass}: the pre-peak compliance reproduced LEFM
beam theory to within $\sim 8\,\%$, but the LEFM peak load
$P_c \approx 6.09~\text{N/mm}$ and post-peak softening could not be
reached on any of four BC variants tried (corner load, distributed load,
displacement vs.\ force, small vs.\ finite strain). This section revisits
the DCB now that Box~4.4 arc-length is operational, to test whether the
prior failure was a Newton-convergence problem (which AL would fix) or a
structural / formulation issue (which AL cannot fix).

%=====================================================
\subsection{Boundary conditions and the role of the MMB fixture (Fig.~8)}

Turon \emph{et al.}~(2006) use the same MMB (mixed-mode bending) fixture
of Crews \& Reeder (Fig.~8 in the paper) for \emph{all} mode-mixity test
cases — DCB, MMB, and ENF — by varying the lever length $c$ which sets
the proportion of end load to middle load applied to the specimen:

\begin{table}[h]
\centering
\small
\begin{tabular}{lll}
\toprule
$G_{II}/G_T$ & test name & loading on the specimen \\
\midrule
$0\,\%$         & \textbf{DCB} (pure mode I)  & end loads only, $\pm P$ on the two arms \\
$20, 50, 80\,\%$ & MMB (mixed)                & end \emph{and} middle loads, ratio set by $c$ \\
$100\,\%$       & ENF (pure mode II)         & middle load only \\
\bottomrule
\end{tabular}
\end{table}

For pure DCB, the lever degenerates: only the end loads are non-zero,
the middle load is zero, and the load-point displacement of the lever
reduces to the relative end opening of the specimen,
$\delta = u_y\big|_{\rm top,\,x=0} - u_y\big|_{\rm bottom,\,x=0}$.
The fixture is thus not actually used; pure DCB is straightforward
end-loaded mode I.

Our setup applies $+\lambda\,\sigma_{\rm ref}$ on \texttt{left\_top}
and $-\lambda\,\sigma_{\rm ref}$ on \texttt{left\_bot}, with
$\sigma_{\rm ref} = P_c^{\rm LEFM}/h_{\rm arm}$, and reports
\texttt{opening = tip\_disp\_top - tip\_disp\_bot}. This is exactly the
DCB ($G_{II}/G_T = 0\,\%$) configuration of Turon's MMB fixture, so the
boundary conditions match the paper.

The plane-strain LEFM peak per unit out-of-plane width is
\begin{equation}
P_c \;=\; \sqrt{\frac{E_{\rm eff}\,h^3\,G_{Ic}}{12\,a_0^2}}
\;=\; \sqrt{\frac{(E/(1-\nu^2))\,h^3\,G_{Ic}}{12\,a_0^2}}
\;=\; 6.09~\text{N/mm},
\end{equation}
which over Turon's $25.4$~mm width is $154.6$~N — matching his
predicted $152.4$~N within $1.4\,\%$. (Earlier sections used the
plane-stress formula $5.89$~N/mm; the plane-strain factor
$1/(1-\nu^2) = 1.067$ has been retained throughout this revision.)

%=====================================================
\subsection{Setup}

Same DCB geometry, mesh, and material as
\texttt{tier2\_1\_dcb\_force.i}:
\begin{itemize}
\item Geometry: $L = 102$~mm, $h = 1.56$~mm per arm,
      $a_0 = 32.9$~mm, plane strain, unit out-of-plane width.
\item Mesh: $408 \times 8$ structured, $0.25$~mm element size in $x$
      ($\sim 5$ elements per Hillerborg cohesive zone length
      $l_{\rm cz} = 1.34$~mm).
\item Material (AS4/PEEK): $E_1 = 122.7$~GPa, $\nu = 0.25$,
      $N = 80$~MPa, $G_{Ic} = 0.969$~N/mm, $K = 10^{6}$~MPa/mm.
\end{itemize}

The only changes vs.\ the prior force-controlled run:
\begin{enumerate}
\item Two \texttt{ALCoupledScaledNeumannBC}s (top arm $+$,
      bottom arm $-$) replace the \texttt{FunctionNeumannBC}s. Reference
      pressure $\sigma_{\rm ref} = P_c^{\rm LEFM}/h = 3.78$~MPa, so
      $\lambda = 1$ corresponds approximately to the LEFM peak.
\item \texttt{ArcLengthTransient} executioner with $\Delta s = 0.5$ in
      displacement L2-norm, $120$ outer steps.
\item Tiny initial offset $u_y\big|_{\rm upper}^{t=0} = 10^{-7}$~mm to
      bypass the bilinear-law zero-tangent corner at $\delta_n = 0$.
\end{enumerate}
File: \texttt{tier2\_1\_dcb\_AL.i}; verification:
\texttt{verify\_dcb\_AL.py}.

%=====================================================
\subsection{Result}

The arc-length solver converges in $1$--$2$ Newton iterations per outer
step throughout the $120$-step run. The trajectory is essentially
\emph{linear} in $(u_{\rm top},\,P)$:
\begin{equation}
\frac{P}{u_{\rm top}}\Big|_{\rm AL} \;=\; 3.219~\text{N/mm/mm},
\qquad
\frac{P}{u_{\rm top}}\Big|_{\rm LEFM\ beam} \;=\; \frac{E_{\rm eff}\,h^3}{4\,a_0^3}
\;=\; 3.488~\text{N/mm/mm},
\end{equation}
a $7.7\,\%$ systematic deviation matching the prior force-controlled
result, attributable to the cohesive-zone added compliance.

\begin{figure}[h]
\centering
\includegraphics[width=\textwidth]{tier2_1_dcb_AL_compare.png}
\caption{Box~4.4 arc-length DCB vs.\ LEFM beam and the prior
force-controlled run.
\textbf{Left}: load--deflection trace. AL (red) and the prior
force-controlled run (green) sit on the same line through
$P_c^{\rm LEFM}$ and continue past it without inflection.
\textbf{Right}: cohesive-zone activation. The maximum normal traction
along the bond saturates at $N = 80$~MPa (cohesive zone reaches
softening locally), but the maximum normal jump $\delta_n$ stays
$\sim 1\,\%$ of $\delta_f$ — the leading edge does not approach
$d = 1$, the crack does not propagate, and no global limit point
appears.}
\label{fig:dcb_al_compare}
\end{figure}

\paragraph{Linearity check.} Slope on the first $30$ converged AL steps
is $3.2195$; on the last $30$ it is $3.2185$ — ratio $0.9997$.
The trajectory is linear within $0.03\,\%$ over the entire $120$-step
run, well past $P_c^{\rm LEFM}$.

\paragraph{Cohesive activation.} max\,$t_n = 79.6$~MPa is reached at
$\lambda \approx 1.7$ ($P \approx 10.0$~N/mm), i.e.\ at $\sim 1.7\times$
the LEFM peak. The leading-edge $\delta_n$ at the end of the run is
$2.12\times 10^{-4}$~mm = $0.88\,\%$ of $\delta_f$. The cohesive zone is
in softening at the leading edge but is far from full decohesion.

%=====================================================
\subsection{Mesh convergence study: the discrete formulation does not converge}
\label{subsec:mesh_conv}

A four-level convergence study under fixed material parameters and BCs:

\begin{table}[h]
\centering
\small
\begin{tabular}{llccc}
\toprule
mesh & element & $h$ [mm] & $\lambda$ at $\max t_n = N$ & $P$ at saturation [N/mm] \\
\midrule
$nx=102$  & QUAD4 & $1.000$  & never ($\max t_n = 10.5$ MPa at $P = 80$) & — \\
$nx=408$  & QUAD4 & $0.250$  & $1.71$ & $10.1$ \\
$nx=816$  & QUAD4 & $0.125$  & $0.57$ & $3.4$ \\
$nx=1632$ & QUAD4 & $0.0625$ & $0.37$ & $2.2$ \\
$nx=408$  & QUAD8 & $0.250$  & $0.91$ & $5.4$ \\
\midrule
LEFM beam-theory peak: & & & & $P_c = 6.09$ \\
\bottomrule
\end{tabular}
\caption{Cohesive saturation load $P_{\rm sat}$ (the applied load at which
the leading-edge interface traction first reaches the cohesive strength $N$)
as the mesh is refined. Halving $h$ from $0.25 \to 0.125 \to 0.0625$ mm
moves $P_{\rm sat}$ from $10.1 \to 3.4 \to 2.2$ N/mm — not approaching any
limit. Switching to QUAD8 at $h = 0.25$ mm shifts $P_{\rm sat}$ to $5.4$
N/mm (closer to LEFM) but does not change the divergent trend.}
\end{table}

The pre-peak compliance, by contrast, IS mesh-convergent: $P/u_{\rm top}$
sits at $3.20$–$3.69$ across all five runs (LEFM beam-theory $3.49$).
What does \emph{not} converge is the load at which the cohesive zone
becomes active.

\paragraph{Why this happens.}
At the bonded$\leftrightarrow$unbonded crack tip, the bulk linear-elasticity
solution has a $1/\sqrt{r}$ stress singularity. The cohesive zone is
supposed to \emph{regularise} this singularity: a true cohesive element
with its own shape functions integrates the bounded traction $t \le N$
over an element length, distributing the cohesive force across a finite
process zone of length $\ell_{\rm cz}$ set by the constitutive law.

The MOOSE \texttt{CZMInterfaceKernel} evaluates $t = K \cdot \jump{u}$ at
side quadrature points, where $\jump{u}$ is interpolated from the bulk's
side trace. As $h \to 0$, quadrature points sit closer to the singular
tip and report ever-higher $t$. The cohesive law clips at $N$, but the
clipping happens at lower and lower applied load. The cohesive zone
never spreads — it concentrates onto a single quadrature point.

This is the discrete signature of a formulation that \emph{inherits} the
bulk singularity instead of \emph{regularising} it.

\paragraph{What this means for adding zero-thickness elements.}
The convergence study justifies the framework PR. Going to a true
cohesive element type with shape-function continuity to the bulk would:
\begin{enumerate}
\item integrate the cohesive law over an element-natural length scale
      rather than at a single quadrature point;
\item produce a mesh-convergent cohesive activation load (independent of
      $h$ once $h \ll \ell_{\rm cz}$);
\item allow the process zone to spread over multiple elements as load
      increases, which is the precondition for a global limit point.
\end{enumerate}
Mesh refinement and higher-order shape functions in the existing
formulation cannot substitute for this — the convergence study shows
both make the activation problem worse, not better, by sharpening the
inherited singularity rather than regularising it.

\paragraph{Effort estimate for the PR.}
\begin{itemize}
\item New \texttt{libMesh::Elem} subclass for a 2D 4-node and 3D 8-node
      zero-thickness cohesive element: $\sim$500--1000 lines + tests.
      libMesh has no existing element type close enough to derive from.
\item Mesh generator that creates these elements at block-pair
      interfaces (essentially \texttt{BreakMeshByBlockGenerator} +
      element creation): $\sim$200--400 lines.
\item Special quadrature rule for zero-thickness integration.
\item Adapt \texttt{BiLinearMixedModeTraction} and friends to the new
      element type — most of the existing CZM material code reuses with
      minor adjustments: $\sim$200--500 lines.
\item Exodus I/O handler (zero-thickness elements are typically written
      as degenerate prisms).
\item Comprehensive testing on existing CZM benchmarks.
\end{itemize}
Realistic estimate: \textbf{a few weeks to a few months} of focused work
for a developer fluent in both libMesh and MOOSE internals.

\begin{figure}[h]
\centering
\includegraphics[width=\textwidth]{tier2_1_dcb_AL_meshconv.png}
\caption{Mesh convergence study. \textbf{Left}: load--deflection
trajectories. All five configurations sit on essentially the same
linear branch (slope $3.20$--$3.69$ vs LEFM beam $3.49$), confirming
that pre-peak compliance is mesh-convergent. \textbf{Right}:
leading-edge interface traction along the bond. The saturation load
shifts toward zero with mesh refinement (red $\to$ cyan $\to$ olive),
confirming the cohesive activation is \emph{not} mesh-convergent in
this formulation. QUAD8 (orange) is closer to the LEFM peak but still
follows the same divergent pattern.}
\label{fig:dcb_meshconv}
\end{figure}

%=====================================================
\subsection{Why is there no limit point? — diagnosis}
\label{subsec:no_limit_point_diagnosis}

A DCB load--deflection curve has a limit point when the cohesive zone's
contribution to the global compliance grows faster than the bulk's
linear loading. Two mechanisms can do this:
\begin{enumerate}
\item[(i)] \textbf{Process zone spread.} More cohesive material enters
  softening as load increases.
\item[(ii)] \textbf{Effective crack advance.} The leading-edge cohesive
  element reaches $d=1$ ($\delta_n = \delta_f$), the effective crack
  tip jumps forward, and $C_{\rm arm} \propto a_{\rm eff}^3$ grows.
\end{enumerate}
Neither happens at physically meaningful loads in this problem. The
following sub-sections quantify why.

\paragraph{(a) The elastic process zone is too short relative to the mesh.}
With cohesive penalty $K = 10^{6}$~MPa/mm and arm bending stiffness
$EI = E h^3/12 = 38{,}791~\text{N}\!\cdot\!\text{mm}^2$ per unit width, the
beam-on-elastic-foundation decay length is
\begin{equation}
\ell_{\rm bef} \;=\; \!\left(\frac{4\,EI}{K}\right)^{\!1/4}
\;=\; \left(\frac{4 \cdot 38{,}791}{10^{6}}\right)^{\!1/4}
\;\approx\; 0.63~\text{mm}.
\end{equation}
At $0.25$~mm element size, only $\sim 2.5$ elements behind the leading
edge see appreciable elastic cohesive deformation. Hillerborg's
``natural'' cohesive process zone length is much larger:
\begin{equation}
\ell_{\rm cz}^{\rm Hillerborg} \;=\; \frac{9\pi}{32}\,\frac{E\,G_{Ic}}{N^{2}}
\;\approx\; 16~\text{mm},
\end{equation}
i.e.\ $\ell_{\rm cz}/\ell_{\rm bef} \approx 25$. The natural process
zone never has room to develop because the elastic foundation is too
stiff: only the immediate vicinity of the leading edge is even
elastically deformed, the rest of the bond is essentially rigid.
Mechanism (i) is therefore stifled.

\paragraph{(b) The leading edge advances too slowly.}
At each AL step we measure $\max_x \delta_n(x)$ along the bond. From
the $400$-step AL run with $K = 10^{6}$:
\begin{table}[h]
\centering
\small
\begin{tabular}{rrrr}
\toprule
$\lambda$ & $P$ [N/mm] & $\max\delta_n$ [mm] & $\max\delta_n / \delta_f$ \\
\midrule
$1.00$ & $5.89$ ($=\!P_c^{\rm LEFM}$) & $\sim\!\delta_0 = 8\times 10^{-5}$ & $0.33\,\%$ \\
$1.81$ & $10.7$  & $2.12\times 10^{-4}$ & $0.88\,\%$ \\
$3.00$ & $17.7$  & $\sim\!7\times 10^{-4}$ & $\sim\!2.9\,\%$ \\
$6.05$ & $35.7$  & $1.31\times 10^{-3}$ & $5.4\,\%$ \\
\bottomrule
\end{tabular}
\end{table}
The trajectory is linear ($P/u_{\rm top} = 3.22$~N/mm/mm to $0.03\,\%$
across all $400$ steps), and $\max\delta_n$ grows roughly linearly with
$\lambda$ at a rate of $dx_{\delta_n}/d\lambda \approx 4\times 10^{-6}$
mm per unit $\lambda$. Reaching $\delta_n = \delta_f$ at the leading
edge — the prerequisite for crack advance — requires roughly
\begin{equation}
\lambda \;\sim\; \frac{\delta_f}{4\times 10^{-6}}
\;\approx\; 6\,000,
\end{equation}
i.e.\ $P \approx 35{,}000$~N/mm — $6{,}000\!\times$ the LEFM peak.
This is meaningless physically (well past the small-strain validity
window, well past any reasonable interpretation of the model). So
mechanism (ii) is also unreachable in practice.

\paragraph{(c) Test: lowering $K$ doesn't help — it just shifts the bottleneck.}
A companion run (\texttt{tier2\_1\_dcb\_AL\_K1e4.i}) with $K = 10^{4}$
(softer cohesive penalty) gives $\ell_{\rm bef} \approx 1.98$~mm — a
better-resolved process zone — but the cohesive softening onset moves
to $\delta_0 = N/K = 8\times 10^{-3}$~mm (100$\times$ larger). At
$\lambda = 2.97$ ($P = 17.5$~N/mm) the leading edge has only $\delta_n
= 2.45\times 10^{-4}$~mm — still in the elastic penalty branch, with
$\max t_n = 2.45$~MPa, i.e.\ $3\,\%$ of $N$. Softening hasn't even
begun at this load level. The bottleneck simply moves from
``leading-edge advance is slow'' (high $K$) to ``loading required to
reach softening is high'' (low $K$). There's no $K$ for which the
cohesive zone produces a clean limit point on this DCB geometry.

\paragraph{(d) Why this differs from Turon's UEL.}
A user-written zero-thickness cohesive element with shape-function
continuity to the bulk concentrates stress at the leading edge as a
near-singular field. Locally the stress \emph{exceeds} $N$ in a
small neighbourhood, drives $\delta_n \to \delta_f$ on a single
quadrature point of the leading element, and that element fully
decoheres. The effective crack tip advances; the next element starts
its own softening from the elevated $\delta_n$ inherited via shape
continuity; the process repeats. This produces a smooth limit point
at the LEFM peak and the post-peak shape of Turon Fig.~9.

The MOOSE \texttt{BreakMeshByBlockGenerator} $+$
\texttt{CZMInterfaceKernel} formulation, by contrast, samples
$\max t_n$ along the bond and \emph{caps} it at $N$ as soon as any
quadrature point reaches $\delta_n = \delta_0$. There is no
shape-function continuity transmitting stress concentration into the
interior. Each element has to be ``loaded up'' independently to reach
$\delta_f$, which is why the rate $d(\max\delta_n)/d\lambda$ stays
small.

\paragraph{Bottom line.} The MOOSE DCB \emph{does} approach the LEFM
peak load (it crosses $P_c = 5.89$~N/mm at $\lambda = 1.0$, with
max traction $79.6$~MPa), but does not have a limit point there
because no single cohesive element actually fails. Reaching the
post-peak ``Fig.~9'' branch in this formulation would require either
mesh refinement of the order $\Delta x \ll \ell_{\rm bef}$ (so the
leading-edge stress concentration is well-resolved) or a true
cohesive-element formulation. Arc-length is orthogonal to this issue
and does not move the diagnosis.

%=====================================================
\subsection{Conclusion: AL doesn't fix this}

The prior diagnosis stands. The MOOSE \texttt{BreakMeshByBlockGenerator}
$+$ \texttt{CZMInterfaceKernel} formulation produces a cohesive process
zone whose stresses average over each interface segment, in contrast to
Turon's user-written cohesive elements which integrate via shape-function
continuity with the bulk and concentrate stress at the leading edge.
Effectively, MOOSE's CZM is ``softer'' at the crack tip — it spreads the
process zone but doesn't push any single point to $d = 1$ at LEFM-peak
loading.

Box~4.4 arc-length is a numerical \emph{path-following} tool: it traces
the equilibrium curve and handles limit points ($du/dP = 0$ or
snap-back, $du/dP < 0$) where standard Newton fails. If the equilibrium
curve has no limit point, AL just retraces what fixed-load Newton already
gives. Here, the curve is linear all the way through $\lambda \sim 1.7$,
so AL contributes nothing the force-controlled run does not.

What \emph{would} reproduce Turon's Fig.~9 in MOOSE remains the same as
the prior recommendation:
\begin{enumerate}
\item Switch to a true cohesive-element formulation (zero-thickness
      interface element with shape-function continuity to the bulk) —
      significant new code in MOOSE.
\item Switch to a different benchmark whose post-peak shape transfers
      to the BreakMesh $+$ InterfaceKernel formulation.
\item Use ENF (Tier 2.2): unstable shear-mode propagation produces
      snap-back even in MOOSE's formulation, where AL becomes
      load-bearing rather than redundant.
\end{enumerate}

%=====================================================
\subsection{What the AL revisit does verify}

\begin{itemize}
\item AL is operational on a structurally meaningful CZM problem
      ($\sim 6500$ elements, mode-I cohesive law, mixed boundary
      conditions).
\item AL converges in $1$--$2$ Newton iterations per outer step at
      machine-precision residual, all $120$ steps.
\item AL traces the same load--deflection trajectory as fixed-force
      Newton on this problem, confirming that the previously observed
      ``no limit point'' diagnosis was correct: it is not a
      Newton-convergence pathology that AL would have hidden.
\item The bilinear traction--separation law's zero-normal-tangent corner
      at $\delta_n = 0$ requires the same $10^{-7}$~mm initial offset
      workaround used in the snap-back demo
      (Section~\ref{sec:snap_back_czm}).
\end{itemize}

The DCB Tier 2.1 status is unchanged: \emph{pre-peak compliance verified
to within $8\,\%$ of LEFM; full Fig.~9 reproduction remains formulation-limited.}
\end{document}
